% !TeX root = ../qm_chap3.tex

\begin{frame}{算符的线性：为什么如此重要？（1/2）}
  
  
    \iqbsectiontitle{线性算符的直观理解}
    \begin{iqbitemize}
      \item \textbf{叠加性的物理意义}
      \begin{iqbitemize}
        \item 如果系统可以处于状态$\psi_1$或$\psi_2$
        \item 那么也可以处于它们的\textbf{线性组合}（linear combination）$c_1\psi_1 + c_2\psi_2$
        \item 这是\textbf{量子叠加原理}（quantum superposition principle）的数学表达
      \end{iqbitemize}

      \item \textbf{齐次性的物理意义}
      \begin{iqbitemize}
        \item 将波函数整体放大或缩小$c$倍
        \item 不改变系统的物理状态
        \item 因为概率密度$|\psi|^2$会相应缩放，但归一化条件保持比例
      \end{iqbitemize}

      \item \textbf{为什么量子力学算符必须是线性的？}
      \begin{iqbitemize}
        \item 实验观察：量子态满足叠加原理
        \item 数学要求：薛定谔方程是\textbf{线性微分方程}（linear differential equation）
        \item 物理后果：保证概率诠释的一致性
      \end{iqbitemize}
    \end{iqbitemize}
  
\end{frame}

\begin{frame}{算符的线性：为什么如此重要？（2/2）}
  \iqbfontsizemode{small}
  
    \iqbsectiontitle{线性的数学验证}
    \begin{iqbitemize}
      \item \textbf{验证拉普拉斯算符的线性}
      \[
        \nabla^2 = \dfrac{\partial^2}{\partial x^2} + \dfrac{\partial^2}{\partial y^2} + \dfrac{\partial^2}{\partial z^2}
      \]

      \item \textbf{叠加性验证}
      \begin{align*}
        \nabla^2(f+g) &= \dfrac{\partial^2(f+g)}{\partial x^2} + \cdots \\
                  &= \dfrac{\partial^2 f}{\partial x^2} + \dfrac{\partial^2 g}{\partial x^2} + \cdots \\
                  &= \nabla^2 f + \nabla^2 g
      \end{align*}

      \item \textbf{齐次性验证}
      \begin{align*}
        \nabla^2(cf) &= \dfrac{\partial^2(cf)}{\partial x^2} + \cdots \\
                 &= c\dfrac{\partial^2 f}{\partial x^2} + \cdots \\
                 &= c\nabla^2 f
      \end{align*}

      \item \textbf{非线性的反例}
      \[
        \sqrt{\hat{\ }}f(x) = \sqrt{f(x)}
      \]
      显然$\sqrt{f+g} \neq \sqrt{f} + \sqrt{g}$，不满足线性
    \end{iqbitemize}
  

\end{frame}

\begin{frame}{线性算符（1/2）}
  
  
    \iqbsectiontitle{线性算符的定义}
    \begin{iqbitemize}
    \item 算符$\hat{A}$是线性的，当且仅当满足两个条件
    \end{iqbitemize}
    \iqbsmallsep

    \iqbsectiontitle{条件1：叠加性（additivity）}
    \begin{iqbitemize}
    \item 数学表达：
    \end{iqbitemize}
      \[
        \hat{A}[f(x) + g(x)] = \hat{A}f(x) + \hat{A}g(x)
      \]
    \begin{iqbitemize}
    \item 物理意义：两个函数之和的算符作用等于各函数算符作用之和
    \item 这允许我们将复杂状态分解为简单状态的叠加
    \end{iqbitemize}
    \iqbsmallsep

    \iqbsectiontitle{条件2：齐次性（homogeneity）}
    \begin{iqbitemize}
    \item 数学表达：
    \end{iqbitemize}
      \[
        \hat{A}[c \cdot f(x)] = c \cdot \hat{A}f(x)
      \]
    \begin{iqbitemize}
    \item 比例系数$c$可以为复数，须验证算符对复数倍函数的作用
    \item 线性算符作用后，波函数的范数仍可能改变
    \end{iqbitemize}
  
\end{frame}

\begin{frame}{线性算符保留量子态的叠加结构}
  \iqblayouttwo{
  
  
    \iqbsectiontitle{为什么量子力学需要线性算符？}
    \begin{iqbitemize}
    \item \textbf{叠加原理的数学基础}
    \subitem{量子系统可以处于多个状态的叠加}
    \subitem{例如：$\Psi = c_1\psi_1 + c_2\psi_2$}
    \item \textbf{态的线性演化}
    \subitem{薛定谔方程是线性的}
    \subitem{初始态的线性叠加演化为末态的线性叠加}
    \item \textbf{测量理论的要求}
    \subitem{本征函数构成完备集}
    \subitem{任意态可展开为本征态的叠加}
    \end{iqbitemize}
    \iqbsmallsep

    
  }{
  
  
    \iqbsectiontitle{常见线性算符实例}
    \begin{iqbitemize}
    \item 微分算符：$\hat{D} = \dfrac{\mathrm{d}}{\mathrm{d}x}$
    \item Laplace算符：$\nabla^2 = \dfrac{\partial^2}{\partial x^2} + \dfrac{\partial^2}{\partial y^2} + \dfrac{\partial^2}{\partial z^2}$
    \item 哈密顿算符：$\hat{H} = -\dfrac{\hbar^2}{2m}\nabla^2 + V$
    \item 位置表象中的乘法算符包括$\hat{x}$和$\hat{V}(x)$，$\hat{p}_x$则是微分算符
    \end{iqbitemize}
    \iqbsmallsep

    \iqborangebox{关键洞察}{态的叠加依赖线性。由自伴哈密顿算符生成的演化还具有幺正性，因而保持总概率。}
  
  }
\end{frame}

\begin{frame}{算符代数：乘法顺序至关重要}
  
  \iqblayouttwo{
    \iqbsectiontitle{算符的加减}
    \begin{iqbitemize}
      \item 两个算符之和定义为：
      \[
        (\hat{A} + \hat{B})f(x) = \hat{A}f(x) + \hat{B}f(x)
      \]
      \item 两个算符之差定义为：
      \[
        (\hat{A} - \hat{B})f(x) = \hat{A}f(x) - \hat{B}f(x)
      \]
      \item 算符加减后仍是算符
      \item 加法满足交换律：$\hat{A} + \hat{B} = \hat{B} + \hat{A}$
    \end{iqbitemize}
  }{
    \iqbsectiontitle{算符的乘积}
    \begin{iqbitemize}
      \item 算符乘积定义为连续作用：
      \[
        \hat{A}\hat{B}f(x) = \hat{A}[\hat{B}f(x)]
      \]
      \item 先作用右边的算符$\hat{B}$，再作用左边的算符$\hat{A}$

      \iqbtinysep

      \item \textbf{重要}：算符乘法一般\textbf{不满足交换律}
      \[
        \hat{A}\hat{B} \neq \hat{B}\hat{A}
      \]

      \iqbtinysep

      \item 这是量子力学与经典力学的重要区别
    \end{iqbitemize}
  }
\end{frame}

\begin{frame}{教材3.1　两种分配律的条件不同}
  \iqblayouttwo{
    \iqbsectiontitle{算符和位于乘积左侧}
    \begin{iqbitemize}
      \item 令任意输入$f$先经过$\hat C$，再应用算符和的定义。
      \[
        \begin{aligned}
          \left[(\hat A+\hat B)\hat C\right]f&=(\hat A+\hat B)(\hat Cf)\\
             &=\hat A(\hat Cf)+\hat B(\hat Cf)\\
             &=(\hat A\hat C+\hat B\hat C)f .
        \end{aligned}
      \]
      \item 对共同定义域中的任意$f$都成立，因此两侧算符相等。推导只用了求和与复合的定义。
    \end{iqbitemize}
  }{
    \iqbsectiontitle{算符和位于乘积右侧}
    \begin{iqbitemize}
      \item 若$\hat A$满足叠加性，可将它分别作用于$\hat Bf$和$\hat Cf$。
      \[
        \begin{aligned}
          \left[\hat A(\hat B+\hat C)\right]f&=\hat A(\hat Bf+\hat Cf)\\
             &=\hat A\hat Bf+\hat A\hat Cf .
        \end{aligned}
      \]
      \item 平方操作$\hat Sf=f^2$作用于$f+g$时额外产生交叉项$2fg$。
      \item 复共轭保留加法，但把常数$c$变为$c^*$。施行两次复共轭得到恒等算符。
    \end{iqbitemize}
  }
\end{frame}

\begin{frame}{算符乘积按从右到左的顺序作用}
  
    \iqbsectiontitle{问题设定}
    \begin{iqbitemize}
    \item 设$\hat{A} = \dfrac{\mathrm{d}}{\mathrm{d}x}$（微分算符），$\hat{B} = x$（乘法算符）
    \item 计算$\hat{A}\hat{B}$和$\hat{B}\hat{A}$并比较
    \item 验证它们是否相等（即是否对易）
    \end{iqbitemize}
    \iqbsmallsep

    \iqbsectiontitle{计算$\hat{A}\hat{B}f(x)$}
    \begin{iqbitemize}
    \item 先乘法后求导：
    \end{iqbitemize}
      \[
        \hat{A}\hat{B}f(x) = \dfrac{\mathrm{d}}{\mathrm{d}x}[x f(x)]
      \]
    \begin{iqbitemize}
    \item 使用乘积法则：$(uv)' = u'v + uv'$
    \end{iqbitemize}
      \[
        = f(x) + x f'(x) = (1 + x\dfrac{\mathrm{d}}{\mathrm{d}x})f(x)
      \]
  
\end{frame}

\begin{frame}{调换求导与乘法的顺序会多出恒等算符}
  \iqbfontsizemode{small}
  \iqblayouttwo{
    \iqbsectiontitle{计算$\hat{B}\hat{A}f(x)$}
    \begin{iqbitemize}
    \item 先求导后乘法：
    \end{iqbitemize}
      \[
        \hat{B}\hat{A}f(x) = x \cdot \dfrac{\mathrm{d}}{\mathrm{d}x}f(x)
      \]
    \begin{iqbitemize}
    \item 直接对$f(x)$求导再乘$x$：
    \end{iqbitemize}
      \[
        = x f'(x) = (x\dfrac{\mathrm{d}}{\mathrm{d}x})f(x)
      \]
  }{
    \iqbsectiontitle{对比与结论}
    \begin{iqbitemize}
    \item $\hat{A}\hat{B} = 1 + x\dfrac{\mathrm{d}}{\mathrm{d}x}$
    \item $\hat{B}\hat{A} = x\dfrac{\mathrm{d}}{\mathrm{d}x}$
    \item \textbf{相差恒等算符}：
    \end{iqbitemize}
      \[
        \hat{A}\hat{B} - \hat{B}\hat{A} = \hat{1} \quad \Rightarrow \quad [\hat{A}, \hat{B}] = \hat{1}
      \]
    \begin{iqbitemize}
    \item 这就是著名的对易关系！
    \end{iqbitemize}
    \iqbsmallsep

    \iqbgreenbox{物理意义}{算符顺序改变结果，是量子力学区别于经典力学的本质特征}
  }
\end{frame}

\begin{frame}{例题3.1：验证算符代数恒等式}
  \iqbfontsizemode{small}
  \iqbsectiontitle{验证$(\hat{A} + \hat{B})^2 \neq \hat{A}^2 + 2\hat{A}\hat{B} + \hat{B}^2$}

  \textbf{问题}：设$\hat{A} = \dfrac{\mathrm{d}}{\mathrm{d}x}$，$\hat{B} = x$，验证算符平方展开

  \iqbsep

  \begin{iqbitemize}
    \item \textbf{经典代数}：$(a+b)^2 = a^2 + 2ab + b^2$

    \item \textbf{算符代数}：一般\textcolor{red}{不成立}！
    \[
      (\hat{A} + \hat{B})^2 = (\hat{A} + \hat{B})(\hat{A} + \hat{B}) = \hat{A}^2 + \hat{A}\hat{B} + \hat{B}\hat{A} + \hat{B}^2
    \]

    \item \textbf{验证}：左边$(\hat{D} + \hat{x})^2 = (\hat{D} + \hat{x})(\hat{D} + \hat{x})$
    \begin{align*}
      &= \hat{D}^2 + \hat{D}\hat{x} + \hat{x}\hat{D} + \hat{x}^2 \\
      &= \hat{D}^2 + (\hat{1} + \hat{x}\hat{D}) + \hat{x}\hat{D} + \hat{x}^2 \\
      &= \hat{D}^2 + \hat{1} + 2\hat{x}\hat{D} + \hat{x}^2
    \end{align*}

    \item \textbf{右边}：$\hat{A}^2 + 2\hat{A}\hat{B} + \hat{B}^2 = \hat{D}^2 + 2\hat{D}\hat{x} + \hat{x}^2$

    \item \textbf{比较}：
    \begin{iqbitemize}
      \item 左边展开：$\hat{D}^2 + \hat{1} + 2\hat{x}\hat{D} + \hat{x}^2$
      \item 普通公式给出的交叉项为$2\hat{D}\hat{x}=2\hat{1}+2\hat{x}\hat{D}$
      \item 正确展开减去普通公式的结果，差值为$-\hat{1}$
    \end{iqbitemize}
  \end{iqbitemize}

  \iqbgreenbox{结论}{算符代数必须严格按照定义展开，不能套用经典代数公式}
\end{frame}

\begin{frame}{例题3.2：算符多项式的计算}
  \iqbfontsizemode{small}
  \iqbsectiontitle{计算$(\hat{A}+\hat{B})^3$的展开式}

  \textbf{问题}：展开算符多项式$(\hat{A}+\hat{B})^3$并说明需要注意的问题

  \iqbmicrosep

  \begin{iqbitemize}
    \item \textbf{错误的展开}（套用经典代数）：
    \[
      (\hat{A}+\hat{B})^3 \neq \hat{A}^3 + 3\hat{A}^2\hat{B} + 3\hat{A}\hat{B}^2 + \hat{B}^3
    \]

    \item \textbf{正确的展开}（保持算符顺序）：
    \begin{align*}
      (\hat{A}+\hat{B})^3 &= (\hat{A}+\hat{B})(\hat{A}+\hat{B})(\hat{A}+\hat{B}) \\
      &= (\hat{A}+\hat{B})(\hat{A}^2 + \hat{A}\hat{B} + \hat{B}\hat{A} + \hat{B}^2) \\
      &= \hat{A}^3 + \hat{A}^2\hat{B} + \hat{A}\hat{B}\hat{A} + \hat{A}\hat{B}^2 \\
      &\quad + \hat{B}\hat{A}^2 + \hat{B}\hat{A}\hat{B} + \hat{B}^2\hat{A} + \hat{B}^3
    \end{align*}

    \item \textbf{关键区别}：
    \begin{iqbitemize}
      \item 普通数相乘可交换，$a^2b=aba=ba^2$，所以这三项可合并为$3a^2b$。
      \item 算符乘积中，$\hat{A}^2\hat{B}$、$\hat{A}\hat{B}\hat{A}$和$\hat{B}\hat{A}^2$一般各不相同。
      \item 一般展开式保留8项及其各自的作用顺序。
    \end{iqbitemize}

    \item \textbf{特例}：如果$\hat{A}$和$\hat{B}$对易（即$[\hat{A},\hat{B}]=0$），则经典公式成立
    \[
      [\hat{A},\hat{B}]=0 \Rightarrow \hat{A}\hat{B} = \hat{B}\hat{A}
    \]
    此时8项可以合并为4项
  \end{iqbitemize}

  \iqbgreenbox{关键}{算符乘法必须保持顺序，只有对易的算符才能交换}
\end{frame}

\begin{frame}{算符的幂运算}
  
  \iqblayouttwo{
    \iqbsectiontitle{定义与性质}
    \begin{iqbitemize}
      \item 算符的$n$次幂定义为连续作用$n$次
      \[
        \hat{A}^n f(x) = \underbrace{\hat{A}\hat{A}\cdots\hat{A}}_{n\text{次}}f(x)
      \]

      \item 例如：$(\dfrac{\mathrm{d}}{\mathrm{d}x})^2 f(x) = \dfrac{\mathrm{d}^2f}{\mathrm{d}x^2}$

      \item 指数律成立：$\hat{A}^m \hat{A}^n = \hat{A}^{m+n}$

      \item 算符函数可用幂级数定义
      \[
        e^{\hat{A}} = \hat{1} + \hat{A} + \dfrac{\hat{A}^2}{2!} + \dfrac{\hat{A}^3}{3!} + \cdots
      \]
    \end{iqbitemize}
  }{
    \iqbsectiontitle{实际应用}
    \begin{iqbitemize}
      \item \textbf{动能算符}：
      \[
        \hat{T} = -\dfrac{\hbar^2}{2m}\dfrac{\mathrm{d}^2}{\mathrm{d}x^2}
      \]
      涉及二阶导数算符

      \item \textbf{哈密顿算符}：
      \[
        \hat{H} = \hat{T} + \hat{V}
      \]
      动能算符与势能算符之和

      \item \textbf{平移算符}：
      \[
        e^{a\dfrac{\mathrm{d}}{\mathrm{d}x}}f(x) = f(x+a)
      \]
      指数算符产生函数平移
    \end{iqbitemize}

    \iqbtinysep

    \iqborangebox{注意}{算符函数的级数展开必须收敛才有意义}
  }
\end{frame}

\begin{frame}{算符多项式与函数（1/2）}
  \iqbfontsizemode{small}
  
    \iqbsectiontitle{算符多项式的定义}
    \begin{iqbitemize}
    \item 算符多项式是算符幂的线性组合
    \end{iqbitemize}
      \[
        P(\hat{A}) = c_0\hat{1} + c_1\hat{A} + c_2\hat{A}^2 + \cdots + c_n\hat{A}^n
      \]
    \begin{iqbitemize}
    \item \textbf{单算符多项式}：类似普通代数
    \[
      (\hat{A} + \hat{1})^2 = \hat{A}^2 + 2\hat{A} + \hat{1}
    \]
    \item 因为算符与自身对易：$[\hat{A}, \hat{A}] = 0$
    \end{iqbitemize}
    \iqbsmallsep

    \iqbsectiontitle{双算符多项式的复杂性}
    \begin{iqbitemize}
    \item \textbf{常见错误}：
    \[
      (\hat{A} + \hat{B})^2 \neq \hat{A}^2 + 2\hat{A}\hat{B} + \hat{B}^2
    \]
    \item 原因：$\hat{A}$与$\hat{B}$不对易
    \end{iqbitemize}
    \iqbtinysep

    \iqbsectiontitle{正确展开方式}
    \begin{iqbitemize}
    \item 保留算符顺序：
    \[
      (\hat{A} + \hat{B})^2 = \hat{A}^2 + \hat{A}\hat{B} + \hat{B}\hat{A} + \hat{B}^2
    \]
    \item 使用对易子简化：
    \[
      (\hat{A} + \hat{B})^2 = \hat{A}^2 + 2\hat{A}\hat{B} + \hat{B}^2 + [\hat{B}, \hat{A}]
    \]
    \end{iqbitemize}
  
\end{frame}

\begin{frame}{算符多项式与函数（2/2）}
  \iqbfontsizemode{small}
  \iqblayouttwo{
    \iqbsectiontitle{算符函数的定义}
    \begin{iqbitemize}
    \item 通过\textbf{泰勒级数展开}（Taylor series）定义
    \[
      f(\hat{A}) = \sum_{n=0}^{\infty} \dfrac{f^{(n)}(0)}{n!}\hat{A}^n
    \]
    \item \textbf{指数算符}（exponential operator）：
    \end{iqbitemize}
      \[
        e^{\hat{A}} = \hat{1} + \hat{A} + \dfrac{1}{2!}\hat{A}^2 + \dfrac{1}{3!}\hat{A}^3 + \cdots
      \]
    \begin{iqbitemize}
    \item 物理中的重要应用：
    \subitem{时间演化算符：$e^{-i\hat{H}t/\hbar}$}
    \subitem{平移算符：$e^{a\mathrm{d}/\mathrm{d}x}$}
    \subitem{旋转算符：$e^{-i\hat{L}_z\phi/\hbar}$}
    \end{iqbitemize}
  }{
    \iqbsectiontitle{算符函数的性质}
    \begin{iqbitemize}
    \item 级数必须收敛才有意义
    \item 不同算符的函数：$f(\hat{A}, \hat{B})$展开更复杂
    \item 特殊情况：若$[\hat{A}, \hat{B}] = 0$，则$e^{\hat{A}}e^{\hat{B}} = e^{\hat{A}+\hat{B}}$
    \item 一般情况：Baker-Campbell-Hausdorff公式
    \end{iqbitemize}
    \iqbsmallsep

    \iqbgreenbox{物理意义}{算符函数描述量子态随时间、空间或角度的演化}
  }
\end{frame}
